% AI-written, AI cross-reviewed; no human sign-off; not a claim that RH is proved
% Atlas unit D108. Local manuscript v3, item G.116, heading on printed PDF p. 186.
% UpperTheory.V.X.residual_bernoulli_nonvanishing
% paper2/decomposition_tables.tex:1012–1017. Verbatim excerpt; SHA-256 e66e1e8065afaa37563a1cd32d885e342d134874109d0dbd79b2e7bfd34ad069.
% Original TeX uses the macros and cross-reference keys of the source manuscript.

\begin{formalstatement}[D108]\label{decomp:D108}
For an odd prime and a nontrivial character of order prime to it whose reduction differs from cyclotomic, a unit first Bernoulli value for the inverse character forces every residual cocycle satisfying UnrEv to be a coboundary.
\end{formalstatement}
{\small\raggedright\noindent\textsf{\bfseries Related reference}\quad \nolinkurl{UpperTheory.V.X.residual_bernoulli_nonvanishing}\par\noindent \nolinkurl{Development/AutomorphicArithmetic/AutomorphicVarieties/LevelCohomology/ResidualBernoulliNonvanishing.lean:174--263}\par\noindent SHA-256 \nolinkurl{e3df603f0ce361270ec3613008c6df20b46e5ffb2a8db4e5ba0f5df75f842ce0}\par\noindent Dominated 242\quad Disjoint 242\quad terminal.\par}
The application concerns a reduced character, an inverse-character Bernoulli unit, the noncyclotomic requirement, and a simultaneous explicit coboundary at every Galois element.
\hyperref[lit:S028]{Published sources and search S028}.
